\documentclass[11pt]{article}
\usepackage[margin=1in]{geometry}
\usepackage{booktabs}
\usepackage{enumitem}
\usepackage{hyperref}

\title{Current PURE Implementation Status}
\author{Internal Notes}
\date{May 2026}

\begin{document}
\maketitle

\section{Naming and Scope}
The project name used in report-facing text should be \textbf{PURE}: \emph{Predicate-aware Uncropped Relation Embedding}. The companion diagnostic dataset and benchmark should be \textbf{CORE}: \emph{Compositional Relation Evaluation}. Avoid the older labels \texttt{PURE-v3}, \texttt{PURE-SPOA}, and \texttt{COMPREL-OV} in final report sections unless they are explicitly described as historical draft names.

\section{Maintained Architecture}
PURE is a one-stage relation model for VG150-style scene graph prediction. It keeps dense CLIP vision tokens active, grounds object nodes with box-conditioned deformable routing, enriches entities with optional language anchors, builds ordered subject--object relation embeddings, and applies an optional relation-context Transformer for graph-level message passing. Predicate scoring combines a lightweight predicate classifier, CLIP-text predicate scoring, and optional Bayesian frequency calibration.

\section{Training Objective}
Report the objective as three main pillars rather than a long list of independent losses:
\begin{enumerate}[leftmargin=1.5em]
  \item \textbf{PredCE-LA}: predicate cross-entropy with optional logit adjustment for long-tail predicates;
  \item \textbf{Counterfactual-SPOA}: queue-based text alignment with role/predicate counterfactual regularization inside the SPOA alignment term;
  \item \textbf{Dense Grounding}: query-to-pair grounding that keeps visual evidence tied to the dense image field.
\end{enumerate}
Visual hard negatives and bilinear mixing are no longer maintained defaults. They are ablation/debug switches only and should not be presented as core report contributions.

\section{Curriculum}
The maintained code path is a single PURE script controlled by \texttt{PURE\_PHASE} rather than separate base/heavy models:
\begin{itemize}[leftmargin=1.5em]
  \item \texttt{core}: freeze CLIP, disable object language anchors, disable relation context, disable frequency calibration. This phase forces deformable routing and the relation decoder to learn visual evidence.
  \item \texttt{scaling}: resume from core, enable object language anchors and relation-context message passing, and optionally use train-time logit adjustment.
  \item \texttt{eval}: keep the scaling architecture and sweep Bayesian/frequency calibration weights for transparent reporting.
\end{itemize}
This curriculum is the report-friendly alternative to the older branch-ramp debugging workflow.

\section{Branch-Ramp Interpretation}
Branch-ramp runs are diagnostic engineering experiments used to identify which modules help or hurt. They are not the final methodological story. If mentioned in the report, they should appear only as an implementation note or ablation rationale: visual-hard negatives and bilinear mixing were tested, found unstable for mean recall, and moved out of the maintained path.

\section{CORE Status}
CORE is a real generated diagnostic dataset with several thousand images organized into six relation groups through the pipeline described in the report draft. It has not yet been used as the primary training or evaluation source for PURE. The final report should state this clearly and frame CORE as a future evaluation/training resource once its scale, audits, and protocol are sufficiently stable.

\section{Immediate Reporting Guidance}
\begin{itemize}[leftmargin=1.5em]
  \item Use \textbf{PURE} and \textbf{CORE} consistently, with full English names at first mention.
  \item Do not claim SOTA unless the exact protocol, split, graph constraint, object-label source, and calibration setting match the compared methods.
  \item Separate no-prior model metrics from eval-time Bayesian/frequency calibration metrics.
  \item Report SGCls only with predicted object labels, e.g., CLIP object classification; GT object labels are PredCls-only.
  \item Move CORE training/evaluation claims to future work unless new completed runs are available.
\end{itemize}

\end{document}